\section{Arquitectura del Sistema}
\label{sec:s0_arquitectura}

La arquitectura de \textbf{EthosHub} se documenta mediante el \textbf{modelo C4} (Context,
Containers, Components, Code), del que en el Sprint 0 se formalizan los dos primeros niveles:
el \textit{Diagrama de Contexto} (Sección~\ref{ssec:s0_c4_contexto}) y el \textit{Diagrama de
Contenedores} (Sección~\ref{ssec:s0_c4_contenedores}). Esta elección satisface el requisito de
ISO/IEC/IEEE 26515 de describir el sistema en niveles progresivos de detalle, aptos tanto para
\textit{stakeholders} de negocio como para el equipo técnico.\par

\subsection{Diagrama de Contexto (Modelo C4) y Alcance General}
\label{ssec:s0_c4_contexto}

El nivel de contexto sitúa a EthosHub como un sistema central que interactúa con dos actores
humanos (Profesional y Reclutador, además del Administrador) y con un conjunto de sistemas
externos de los que depende: \textbf{Supabase} (autenticación y \textit{Realtime}), una
\textbf{base de datos PostgreSQL 15.10} autoalojada en el servidor TIS de la UMSS
(administrada vía phpPgAdmin), \textbf{Google Gemini} (servicios de IA), \textbf{Google Drive}
(importación de evidencias) y un \textbf{servidor SMTP} (correo transaccional). La
Figura~\ref{fig:s0_c4_contexto} resume estas relaciones.

\begin{figure}[!ht]
    \centering
    \begin{tikzpicture}[
        node distance=1.1cm and 1.4cm,
        actor/.style={draw=colorPrimario, fill=headerTabla, thick, rounded corners=2pt,
            text width=2.3cm, align=center, minimum height=1.0cm, font=\sffamily\footnotesize},
        system/.style={draw=colorPrimario, fill=colorPrimario, text=white, thick,
            rounded corners=3pt, text width=3.2cm, align=center, minimum height=1.4cm,
            font=\sffamily\bfseries\small},
        ext/.style={draw=grisISO, fill=grisTecnico, thick, rounded corners=2pt,
            text width=2.4cm, align=center, minimum height=1.0cm, font=\sffamily\scriptsize},
        rel/.style={-{Latex[length=2mm]}, grisISO, font=\sffamily\tiny}
    ]
        \node[system] (sys) {EthosHub\\\scriptsize Portafolios \& Talento};
        \node[actor, left=of sys] (prof) {Profesional};
        \node[actor, above=of sys] (recr) {Reclutador};
        \node[ext, right=of sys] (supa) {Supabase\\(Auth/DB/Realtime)};
        \node[ext, below left=of sys] (gem) {Google Gemini (IA)};
        \node[ext, below=of sys] (drive) {Google Drive};
        \node[ext, below right=of sys] (smtp) {SMTP (correo)};

        \draw[rel] (prof) -- node[above]{publica portafolio} (sys);
        \draw[rel] (recr) -- node[right]{descubre talento} (sys);
        \draw[rel] (sys) -- node[above]{JWT/SQL/RPC} (supa);
        \draw[rel] (sys) -- node[left]{\textit{insights}} (gem);
        \draw[rel] (sys) -- node[left]{evidencias} (drive);
        \draw[rel] (sys) -- node[right]{OTP/avisos} (smtp);
    \end{tikzpicture}
    \caption{Diagrama de Contexto C4 (Nivel 1) de EthosHub y sus sistemas externos.}
    \label{fig:s0_c4_contexto}
\end{figure}

El \textbf{alcance general} acordado para el producto se segmentó en ocho macro-funcionalidades
(Épicas), que se detallan en la Sección~\ref{ssec:s0_backlog}. La autenticación se
\textbf{delega en Supabase Auth}; el backend actúa como \textit{Resource Server} y únicamente
valida los JWT emitidos, operando sobre el schema \comando{core} con jOOQ.

\subsection{Diagrama de Contenedores y Tecnologías Seleccionadas}
\label{ssec:s0_c4_contenedores}

En el nivel de contenedores se materializa la decisión arquitectónica clave del Sprint 0: el
\textbf{despliegue monolítico}. El \textit{build} de producción del frontend (\comando{dist/})
se copia a \comando{src/main/resources/static/} del backend y se empaqueta dentro de un único
JAR ejecutable de Spring Boot. Un filtro \comando{SpaFallbackFilter} sirve \comando{index.html}
para las rutas de la SPA, mientras que los controladores REST atienden \comando{/api},
\comando{/auth} y \comando{/v1}. La Figura~\ref{fig:s0_c4_contenedores} ilustra el patrón
\textit{Controllers $\to$ Services $\to$ Repositories} (jOOQ).

\begin{figure}[!ht]
    \centering
    \begin{tikzpicture}[
        layer/.style={draw=colorPrimario, fill=headerTabla, thick, rounded corners=2pt,
            text width=3.0cm, align=center, minimum height=0.9cm, font=\sffamily\scriptsize},
        store/.style={draw=grisISO, fill=grisTecnico, thick, rounded corners=2pt,
            text width=3.0cm, align=center, minimum height=0.9cm, font=\sffamily\scriptsize},
        jar/.style={draw=colorPrimario, thick, dashed, rounded corners=4pt, inner sep=10pt},
        rel/.style={-{Latex[length=2mm]}, grisISO, font=\sffamily\tiny}
    ]
        \node[layer] (spa) {SPA React (static/)};
        \node[layer, below=0.7cm of spa] (ctrl) {Controllers (REST)};
        \node[layer, below=0.7cm of ctrl] (svc) {Services (dominio)};
        \node[layer, below=0.7cm of svc] (repo) {Repositories (jOOQ)};
        \node[store, right=2.2cm of repo] (db) {PostgreSQL 15.10\\(servidor TIS)};
        \node[store, right=2.2cm of svc] (apis) {Gemini / SMTP / Drive};

        \draw[rel] (spa) -- node[right]{fetch \comando{/api}} (ctrl);
        \draw[rel] (ctrl) -- (svc);
        \draw[rel] (svc) -- (repo);
        \draw[rel] (repo) -- node[above]{SQL / RPC} (db);
        \draw[rel] (svc) -- node[above]{REST} (apis);

        \begin{scope}[on background layer]
            \node[jar, fit=(spa)(ctrl)(svc)(repo)] (box) {};
        \end{scope}
        \node[font=\sffamily\bfseries\tiny\color{colorPrimario}, above=1pt of box.north]
            {JAR monolítico EthosHub};
    \end{tikzpicture}
    \caption{Diagrama de Contenedores C4 (Nivel 2): empaquetado monolítico FE+BE.}
    \label{fig:s0_c4_contenedores}
\end{figure}

La tabla~\ref{tab:s0_stack} justifica cada tecnología seleccionada. La elección prioriza la
\textbf{seguridad delegada} (Supabase Auth + RLS), el \textbf{SQL tipado} (jOOQ frente a un ORM
de propósito general) y la \textbf{cohesión de despliegue} (un solo artefacto).

\begin{table}[!ht]
    \centering
    \renewcommand{\arraystretch}{1.35}
    \fontsize{9}{10.8}\selectfont\sffamily
    \begin{tabularx}{\textwidth}{|l|l|X|}
        \hline
        \rowcolor{colorPrimario}
        \textcolor{white}{\textbf{Capa}} & \textcolor{white}{\textbf{Tecnología}} & \textcolor{white}{\textbf{Justificación}} \\ \hline
        Lenguaje (BE) & Java 17 & LTS, \textit{records}, \textit{pattern matching} y rendimiento estable. \\ \hline
        \rowcolor{grisTecnico}
        Framework (BE) & Spring Boot 4.0.6 & Web MVC, Security, Validation y Actuator integrados. \\ \hline
        Acceso a datos & jOOQ + Spring JDBC & SQL tipado contra el schema \comando{core}; control total de consultas. \\ \hline
        \rowcolor{grisTecnico}
        Base de datos & PostgreSQL 15.10 (servidor TIS) & RLS, RPCs \comando{SECURITY DEFINER}; autoalojada y administrada vía phpPgAdmin. \\ \hline
        Seguridad & Spring Security (OAuth2 RS) & Validación de JWT de Supabase vía JWKS. \\ \hline
        \rowcolor{grisTecnico}
        Framework (FE) & React 18.3 + TS 5.3 & Componentes, \textit{strict mode} y tipado estático. \\ \hline
        \textit{Bundler} (FE) & Vite 5.1 & HMR veloz y \textit{build} optimizado con \textit{code splitting}. \\ \hline
        \rowcolor{grisTecnico}
        Estado (FE) & Zustand 4.5 & \textit{Stores} por dominio, sin \textit{boilerplate}. \\ \hline
        Documentación API & springdoc-openapi 2.7 & Swagger UI en \comando{/swagger-ui.html}. \\ \hline
    \end{tabularx}
    \caption{Stack tecnológico seleccionado y su justificación.}
    \label{tab:s0_stack}
\end{table}

\begin{NotaTecnica}
La organización del backend sigue un patrón de \textbf{\textit{vertical slice}}: cada dominio
funcional (\comando{auth}, \comando{project}, \comando{recruiter}, etc.) es un paquete que
contiene su \comando{Controller}, \comando{Service}, \comando{Repository} y sus \comando{dto/}.
Este principio, fijado en el Sprint 0, se mantiene invariante a lo largo de todos los sprints.
\end{NotaTecnica}

\subsection{Vista de Componentes (Modelo C4 — Nivel 3)}
\label{ssec:s0_c4_componentes}
Descendiendo al tercer nivel C4, la Figura~\ref{fig:s0_c4_componentes} desglosa los componentes
internos de un \textit{vertical slice} típico y su relación con la seguridad transversal.

\begin{figure}[!ht]
    \centering
    \begin{tikzpicture}[font=\sffamily\scriptsize, >=Latex, node distance=0.7cm and 1.0cm,
        comp/.style={draw=colorPrimario, fill=headerTabla, rounded corners=2pt,
            text width=2.5cm, align=center, minimum height=0.9cm},
        sec/.style={draw=grisISO, fill=grisTecnico, rounded corners=2pt,
            text width=2.5cm, align=center, minimum height=0.9cm},
        rel/.style={-{Latex[length=1.8mm]}, colorPrimario, font=\tiny}]
        \node[comp] (ct) {Controller};
        \node[comp, right=of ct] (sv) {Service};
        \node[comp, right=of sv] (rp) {Repository (jOOQ)};
        \node[sec, below=of sv] (fil) {Filtros de seguridad};
        \node[sec, below=of rp] (db) {core (RLS) / RPC};
        \draw[rel] (ct) -- node[above]{DTO} (sv);
        \draw[rel] (sv) -- node[above]{entidad} (rp);
        \draw[rel] (fil) -- node[left]{JWT} (ct);
        \draw[rel] (rp) -- node[right]{SQL} (db);
    \end{tikzpicture}
    \caption{Vista de componentes C4 (Nivel 3) de un \textit{vertical slice} con seguridad transversal.}
    \label{fig:s0_c4_componentes}
\end{figure}

\clearpage
